% ============================================================
% ECE 2112: Digital Techniques Sessional -- Lab Report 2
% Body of the report (cover page is handled separately)
% ============================================================
\documentclass[12pt,a4paper]{article}

% ---------- Page geometry ----------
\usepackage[a4paper,margin=1in,headheight=15pt]{geometry}

% ---------- Core packages ----------
\usepackage[T1]{fontenc}
\usepackage{lmodern}
\usepackage{microtype}
\usepackage{amsmath,amssymb}
\usepackage{graphicx}
\usepackage{array}
\usepackage{booktabs}
\usepackage{caption}
\usepackage{float}
\usepackage[colorlinks=true,linkcolor=black,citecolor=black,urlcolor=blue]{hyperref}
\usepackage{fancyhdr}
\usepackage{setspace}
\usepackage{titlesec}
\usepackage{enumitem}
\usepackage{diagbox}

\begin{document}
\thispagestyle{empty}
\begin{center}

{\itshape Heavens Light is Our Guide}

\vspace{0.1cm}

% Uncomment after adding the logo
\includegraphics[width=5cm]{ruet-logo.png}

\vspace{0.2cm}

{\Large\textbf{Rajshahi University of Engineering and Technology}}

\vspace{0.2cm}

{\Large\textbf{Department of Electrical \& Computer Engineering}}

\vspace{1cm}

\textbf{Course Title}

\vspace{0.2cm}

Digital Techniques Sessional

\vspace{0.3cm}

\textbf{Course No: ECE 2112}

\vspace{0.5cm}

\textbf{Lab Report -02}

\vspace{0.8cm}

\textbf{Name of the Experiment: Simplification of Boolean Expressions Using Boolean Algebra and Karnaugh Map, and Verification Using Logic Circuit Simulation}

\vspace{0.5cm}


\end{center}

\vspace{2cm}

\noindent
\begin{tabular}{p{0.45\textwidth} p{0.45\textwidth}}
\textbf{Submitted By:} & \textbf{Submitted To:} \\[0.7cm]
Uddipon Chowdhury & Mst. Mazeda Noor Tasnim \\
Roll: 2410019 & Lecturer \\
Registration No: 1072 & ECE, RUET \\
\end{tabular}
\clearpage

% ---------- Justified body text (LaTeX default) ----------
\setlength{\parindent}{0pt}
\setlength{\parskip}{0.6em}
\onehalfspacing
\captionsetup{position=below,font=small,labelfont=bf}

% ---------- Header / footer ----------
\pagestyle{fancy}
\fancyhf{}
\fancyhead[L]{ECE 2112: DT Sessional}
\fancyhead[R]{Lab Report 2}
\fancyfoot[C]{\thepage}
\renewcommand{\headrulewidth}{0.4pt}

% ---------- Section heading style ----------
\titleformat{\section}{\normalfont\Large\bfseries}{\thesection}{1em}{}
\titleformat{\subsection}{\normalfont\large\bfseries}{\thesubsection}{1em}{}
\titlespacing*{\subsection}{0pt}{1.1em}{0.5em}

% ---------- Complement / bar notation (used consistently throughout) ----------
\newcommand{\bA}{\overline{A}}
\newcommand{\bB}{\overline{B}}
\newcommand{\bC}{\overline{C}}

%\begin{document}

% The cover page is supplied separately and carries no page number.
% This document continues the numbering from page 2 onward.
\setcounter{page}{2}
\pagestyle{reportstyle}

\vspace{1em}

% ============================================================
\clearpage
\section*{Problem 1}
\addcontentsline{toc}{section}{Problem 1}

\subsection*{Problem Statement}
Simplify the Boolean expression $F(A,B,C) = \bA BC + A\bB C + AB\bC + ABC$ using Boolean algebra and the Karnaugh map, and verify both the original and the simplified circuits through logic simulation.

\subsection*{Objective}
The goal is to reduce a three-variable sum-of-minterms expression to its minimal sum-of-products form, confirm that reduction with a Karnaugh map, and check that the original and simplified circuits behave identically by building both in Logisim.

\subsection*{Algebraic Simplification}
\begin{align*}
F(A,B,C) &= \bA BC + A\bB C + AB\bC + ABC \\
         &= \bA BC + ABC + A\bB C + ABC + AB\bC + ABC \\
         &= BC(\bA + A) + AC(\bB + B) + AB(\bC + C) \\
         &= BC + AC + AB \\
         &= AB + AC + BC
\end{align*}

\subsection*{Truth Table}
\begin{table}[H]
\centering
\renewcommand{\arraystretch}{1.15}
\begin{tabular}{|c|c|c|c|}
\hline
\textbf{$A$} & \textbf{$B$} & \textbf{$C$} & \textbf{$F = AB + AC + BC$} \\
\hline
0 & 0 & 0 & 0 \\
\hline
0 & 0 & 1 & 0 \\
\hline
0 & 1 & 0 & 0 \\
\hline
0 & 1 & 1 & 1 \\
\hline
1 & 0 & 0 & 0 \\
\hline
1 & 0 & 1 & 1 \\
\hline
1 & 1 & 0 & 1 \\
\hline
1 & 1 & 1 & 1 \\
\hline
\end{tabular}
\caption{Truth table for $F = AB + AC + BC$}
\label{tab:p1-truth}
\end{table}
\subsection*{Karnaugh Map}
\begin{table}[H]
\centering
\renewcommand{\arraystretch}{1.5}
\begin{tabular}{|c|c|c|c|c|}
\hline
\diagbox{$A$}{$BC$} & \textbf{00} & \textbf{01} & \textbf{11} & \textbf{10} \\
\hline
\textbf{0} & 0 & 0 & 1 & 0 \\
\hline
\textbf{1} & 0 & 1 & 1 & 1 \\
\hline
\end{tabular}
\caption{K-map for $F(A,B,C) = \bA BC + A\bB C + AB\bC + ABC$}
\label{tab:p1-kmap}
\end{table}

Three pairs of adjacent 1-cells can be grouped: the pair at $BC = 11$ (covering both rows) gives $BC$; the pair with $A = 1,\ C = 1$ (i.e.\ $BC = 01, 11$) gives $AC$; and the pair with $A = 1,\ B = 1$ (i.e.\ $BC = 11, 10$) gives $AB$. Combining these three groups produces the minimized expression $F = AB + AC + BC$, which agrees with the algebraic result obtained above.

\subsection*{Circuit Diagram}
\begin{figure}[H]
\centering
\includegraphics[width=\textwidth]{ciruit1.png}
\caption{Original circuit (four 3-input AND gates plus a 3-input OR gate, realizing $\bA BC + A\bB C + AB\bC + ABC$) and simplified circuit (three 2-input AND gates plus an OR gate, realizing $AB + AC + BC$), simulated in Logisim}
\label{fig:p1-circuit}
\end{figure}

\subsection*{Discussion}
The original circuit implements the expression exactly as written: it needs four 3-input AND gates, one for each product term, along with inverters to generate $\bA$, $\bB$, and $\bC$ wherever they appear, and a 3-input OR gate to combine the terms. The simplified circuit, in contrast, realizes $AB + AC + BC$ using only three 2-input AND gates feeding a single OR gate, and it needs no inverters at all. Simulating both versions in Logisim for every possible input combination should produce identical outputs, confirming that the algebraic and Karnaugh-map simplifications are correct. Overall, the simplified circuit is clearly the more efficient design: it uses fewer gates, fewer total gate inputs, and eliminates the inverters entirely.

\subsection*{Conclusion}
The expression $F = \bA BC + A\bB C + AB\bC + ABC$ was successfully reduced to $F = AB + AC + BC$ through Boolean algebra and confirmed independently with a Karnaugh map. Logisim simulation of both the original and the simplified circuits is expected to give identical outputs across all input combinations, verifying that the simplification is correct.

% ============================================================
\clearpage
\section*{Problem 2}
\addcontentsline{toc}{section}{Problem 2}

\subsection*{Problem Statement}
Simplify the Boolean expression $F(A,B,C) = A(A+B)(A+B+C)$ using Boolean algebra and the Karnaugh map, and verify both the original and the simplified circuits through logic simulation.

\subsection*{Objective}
The aim is to reduce a three-variable product-of-sums-type expression to its simplest possible form, confirm that reduction using a Karnaugh map, and check the equivalence of the original and simplified circuits in Logisim.

\subsection*{Algebraic Simplification}
\begin{align*}
F(A,B,C) &= A(A+B)(A+B+C) \\
         &= (A + AB)(A+B+C) \\
         &= A(A+B+C) + AB(A+B+C) \\
         &= A + AB + AC + AB + AB + ABC \\
         &= A + AB + AC + ABC \\
         &= A(1 + B + C + BC) \\
         &= A
\end{align*}

\subsection*{Truth Table}
\begin{table}[H]
\centering
\renewcommand{\arraystretch}{1.15}
\begin{tabular}{|c|c|c|c|}
\hline
\textbf{$A$} & \textbf{$B$} & \textbf{$C$} & \textbf{$F = A$} \\
\hline
0 & 0 & 0 & 0 \\
\hline
0 & 0 & 1 & 0 \\
\hline
0 & 1 & 0 & 0 \\
\hline
0 & 1 & 1 & 0 \\
\hline
1 & 0 & 0 & 1 \\
\hline
1 & 0 & 1 & 1 \\
\hline
1 & 1 & 0 & 1 \\
\hline
1 & 1 & 1 & 1 \\
\hline
\end{tabular}
\caption{Truth table for $F = A$}
\label{tab:p2-truth}
\end{table}

\subsection*{Karnaugh Map}
\begin{table}[H]
\centering
\renewcommand{\arraystretch}{1.5}
\begin{tabular}{|c|c|c|c|c|}
\hline
\diagbox{$A$}{$BC$} & \textbf{00} & \textbf{01} & \textbf{11} & \textbf{10} \\
\hline
\textbf{0} & 0 & 0 & 0 & 0 \\
\hline
\textbf{1} & 1 & 1 & 1 & 1 \\
\hline
\end{tabular}
\caption{K-map for $F(A,B,C) = A(A+B)(A+B+C)$}
\label{tab:p2-kmap}
\end{table}

All four cells in the $A = 1$ row are 1, while the entire $A = 0$ row is 0. These four adjacent cells form a single group spanning the whole row, eliminating both $B$ and $C$ and leaving only $F = A$ -- exactly matching the algebraic simplification above.

\subsection*{Circuit Diagram}
\begin{figure}[H]
\centering
\includegraphics[width=\textwidth]{circuit2.png}
\caption{Original circuit (two OR gates realizing $A+B$ and $A+B+C$, combined with $A$ through an AND gate) and simplified circuit (a direct connection representing $F = A$), simulated in Logisim}
\label{fig:p2-circuit}
\end{figure}

\subsection*{Discussion}
The original circuit first forms $A+B$ with one OR gate, extends that result to $A+B+C$ with a second OR gate, and then ANDs both results with $A$ itself. Since $A$ is common to every term in the expression, all of the higher-order terms are absorbed, and the whole network collapses to the single literal $A$; the simplified circuit is therefore nothing more than a wire connecting input $A$ directly to the output, with no gates required. This is a textbook example of the absorption law $A + AB = A$ being applied twice in succession. Testing both circuits in Logisim for every combination of $A$, $B$, and $C$ should confirm that the output always equals $A$, irrespective of $B$ and $C$.

\subsection*{Conclusion}
The expression $F = A(A+B)(A+B+C)$ was shown to reduce to $F = A$ through repeated application of the absorption law. This result is confirmed by the Karnaugh map and is expected to be validated by Logisim simulation, which should show that $B$ and $C$ have no effect on the output.

% ============================================================
\clearpage
\section*{Problem 3}
\addcontentsline{toc}{section}{Problem 3}

\subsection*{Problem Statement}
Simplify the Boolean expression $F(A,B,C) = \left(A+(BC)'\right)'(AB + ABC)$ using Boolean algebra and the Karnaugh map, and verify both the original and the simplified circuits through logic simulation.

\subsection*{Objective}
The objective is to reduce a mixed AND-OR-complement expression using De Morgan's theorem together with Boolean algebra, confirm the result with a Karnaugh map, and use Logisim to verify that the simplified circuit behaves as a constant.

\subsection*{Algebraic Simplification}
\begin{align*}
F(A,B,C) &= \overline{A + \overline{BC}}\,(AB + ABC) \\
         &= \bA \cdot BC \cdot (AB + ABC) \qquad \text{(De Morgan's theorem)} \\
         &= \bA BC \cdot AB(1+C) \\
         &= \bA BC \cdot AB \\
         &= (A\bA)\,BC\,B \\
         &= 0
\end{align*}

\subsection*{Truth Table}
\begin{table}[H]
\centering
\renewcommand{\arraystretch}{1.15}
\begin{tabular}{|c|c|c|c|}
\hline
\textbf{$A$} & \textbf{$B$} & \textbf{$C$} & \textbf{$F = 0$} \\
\hline
0 & 0 & 0 & 0 \\
\hline
0 & 0 & 1 & 0 \\
\hline
0 & 1 & 0 & 0 \\
\hline
0 & 1 & 1 & 0 \\
\hline
1 & 0 & 0 & 0 \\
\hline
1 & 0 & 1 & 0 \\
\hline
1 & 1 & 0 & 0 \\
\hline
1 & 1 & 1 & 0 \\
\hline
\end{tabular}
\caption{Truth table for $F = 0$}
\label{tab:p3-truth}
\end{table}

\subsection*{Karnaugh Map}
\begin{table}[H]
\centering
\renewcommand{\arraystretch}{1.5}
\begin{tabular}{|c|c|c|c|c|}
\hline
\diagbox{$A$}{$BC$} & \textbf{00} & \textbf{01} & \textbf{11} & \textbf{10} \\
\hline
\textbf{0} & 0 & 0 & 0 & 0 \\
\hline
\textbf{1} & 0 & 0 & 0 & 0 \\
\hline
\end{tabular}
\caption{K-map for $F(A,B,C) = \overline{A + \overline{BC}}\,(AB + ABC)$}
\label{tab:p3-kmap}
\end{table}

Every cell of the map is 0, so there are no 1-cells available to group. A Karnaugh map with no populated cells represents the constant function $F = 0$, which agrees with the algebraic result.

\subsection*{Circuit Diagram}
\begin{figure}[H]
\centering
\includegraphics[width=\textwidth]{circuit3.png}
\caption{Original circuit realizing $\left(A+(BC)'\right)'(AB+ABC)$ using NAND-type and OR/AND stages, and simplified circuit showing the output permanently at logic 0, simulated in Logisim}
\label{fig:p3-circuit}
\end{figure}

\subsection*{Discussion}
Applying De Morgan's theorem to $\overline{A+\overline{BC}}$ converts it into $\bA \cdot BC$; the second factor, $AB + ABC$, then simplifies to $AB$ by absorption, since $AB + ABC = AB(1+C) = AB$. Multiplying $\bA BC$ by $AB$ introduces the term $A \cdot \bA$, which is always 0, so the entire expression is forced to 0 regardless of the values of $A$, $B$, and $C$. Cycling through all eight input combinations on the original circuit in Logisim should show the output remaining at logic 0 throughout, exactly as predicted by the all-zero Karnaugh map.

\subsection*{Conclusion}
The expression $F = \left(A+(BC)'\right)'(AB+ABC)$ was shown to be identically zero for every input combination. This was verified algebraically and confirmed by an empty Karnaugh map; Logisim simulation is expected to validate this by showing a constant logic-0 output regardless of $A$, $B$, and $C$.

% ============================================================
\clearpage
\section*{Problem 4}
\addcontentsline{toc}{section}{Problem 4}

\subsection*{Problem Statement}
Simplify the Boolean expression $F(A,B) = \left[\bB(A+B) + (A+B)(A+\bB)\right]\bB$ using Boolean algebra and the Karnaugh map, and verify both the original and the simplified circuits through logic simulation.

\subsection*{Objective}
The goal is to reduce a two-variable expression built from nested OR and AND terms to its minimal form, confirm the reduction using a Karnaugh map, and verify the equivalence of the original and simplified circuits in Logisim.

\subsection*{Algebraic Simplification}
\begin{align*}
F(A,B) &= \left[\bB(A+B) + (A+B)(A+\bB)\right]\bB \\
       &= (A\bB + B\bB + AA + A\bB + A\bB + B\bB)\bB \\
       &= (A + A\bB + A\bB + A\bB)\bB \\
       &= A(1 + \bB + \bB + \bB)\,\bB \\
       &= A\bB
\end{align*}

\subsection*{Truth Table}
\begin{table}[H]
\centering
\renewcommand{\arraystretch}{1.15}
\begin{tabular}{|c|c|c|}
\hline
\textbf{$A$} & \textbf{$B$} & \textbf{$F = A\bB$} \\
\hline
0 & 0 & 0 \\
\hline
0 & 1 & 0 \\
\hline
1 & 0 & 1 \\
\hline
1 & 1 & 0 \\
\hline
\end{tabular}
\caption{Truth table for $F = A\bB$}
\label{tab:p4-truth}
\end{table}

\subsection*{Karnaugh Map}

\begin{table}[H]
\centering
\renewcommand{\arraystretch}{1.3}
\begin{tabular}{|c|c|c|}
\hline
\diagbox{$A$}{$B$} & \textbf{0} & \textbf{1} \\
\hline
\textbf{0} & 0 & 0 \\
\hline
\textbf{1} & 1 & 0 \\
\hline
\end{tabular}
\caption{K-map for $F(A,B) = \left[\overline{B}(A+B) + (A+B)(A+\overline{B})\right]\overline{B}$}
\label{tab:p4-kmap}
\end{table}
Only a single cell, at $A=1,\ B=0$, is populated. A lone 1-cell cannot be paired with any adjacent cell, so it is read directly as the product of its coordinate literals, giving $F = A\bB$ -- matching the algebraic result.

\subsection*{Circuit Diagram}
\begin{figure}[H]
\centering
\includegraphics[width=\textwidth]{circuit4.png}
\caption{Original circuit realizing $\left[\bB(A+B) + (A+B)(A+\bB)\right]\bB$ using OR, AND, and NOT stages, and simplified circuit reduced to a single AND gate realizing $A\bB$, simulated in Logisim}
\label{fig:p4-circuit}
\end{figure}

\subsection*{Discussion}
The original expression can be simplified by F = \(A\overline{B}\).

Thus, the original multi-gate network is functionally equivalent to a simplified implementation of \(A\overline{B}\), requiring an AND operation with the \(B\) input complemented. Testing both circuits in Logisim for all four combinations of \(A\) and \(B\) should confirm that the output is 1 only when \(A=1\) and \(B=0\), matching the truth table and Karnaugh map.

\subsection*{Conclusion}
The expression $F = \left[\bB(A+B) + (A+B)(A+\bB)\right]\bB$ was successfully reduced to $F = A\bB$ using Boolean algebra and confirmed with a Karnaugh map. Logisim simulation of the original and simplified circuits is expected to show identical outputs across all input combinations.

% ============================================================

\addcontentsline{toc}{section}{References}

\begin{thebibliography}{9}

\bibitem{mano}
M.\ M.\ Mano, \emph{Digital Design}, 5th ed. Upper Saddle River, NJ: Pearson, 2013.

\bibitem{tocci}
R.\ J.\ Tocci, N.\ S.\ Widmer, and G.\ L.\ Moss, \emph{Digital Systems: Principles and Applications}, 11th ed. Upper Saddle River, NJ: Pearson, 2011.

\end{thebibliography}

\end{document}
